%---- Required Packages and Functions ----

\documentclass[a4paper,11pt]{article}
\usepackage{latexsym}
\usepackage{xcolor}
\usepackage{float}
\usepackage{ragged2e}
\usepackage[empty]{fullpage}
\usepackage{wrapfig}
\usepackage{lipsum}
\usepackage{tabularx}
\usepackage{titlesec}
\usepackage{geometry}
\usepackage{marvosym}
\usepackage{verbatim}
\usepackage{enumitem}
\usepackage[hidelinks]{hyperref}
% \usepackage{hyperref}
\usepackage{fancyhdr}
\usepackage{multicol}
\usepackage{multirow}
\usepackage{graphicx}
\usepackage{cfr-lm}
\usepackage[T1]{fontenc}
\usepackage{fontawesome}
\usepackage[most]{tcolorbox}
\usepackage{lastpage}
\usepackage{booktabs}



% Set the margins
\geometry{left=0.85cm, top=0.8cm, right=0.85cm, bottom=0.2cm}


% Sections formatting
\titleformat{\section}{
  \vspace{-4pt}\scshape\raggedright\large
}{}{0em}{}[\color{black}\titlerule \vspace{-7pt}]


%-------------------------
% Custom commands

\newcommand{\resumePOR}[3]{
\vspace{0.5mm}\item[]
    \begin{tabular*}{\textwidth}[t]{l@{\extracolsep{\fill}}r}
    \hspace{-3mm}{#1}:\hspace{1mm} & \hspace*{0pt}\hfill{\footnotesize{ #3}} \vspace{-0.5mm}\\ \hspace{-2.9mm}#2 
    \end{tabular*}
    \vspace{0mm}
}

\newcommand{\resumeExp}[2]{
\vspace{0mm}\item[]
    \begin{tabular*}{\textwidth}[t]{l@{\extracolsep{\fill}}r}
        \hspace{-4.4mm} \small\textbf{#1} & {\footnotesize{#2}} \vspace{-1.2mm} %\\
        %\hspace{-4.3mm} \footnotesize{\text{#2}} & \footnotesize{#4}
    \end{tabular*}
    \vspace*{-7mm}
}

\newcommand{\resumeProject}[2]{
\vspace{0mm}\item[]
    \begin{tabular*}{\textwidth}[t]{l@{\extracolsep{\fill}}r}
        \hspace{-4.4mm} \small\textbf{#1} & {\footnotesize{#2}} \vspace{-1.2mm} %\\
        %\hspace{-4.3mm} \footnotesize{\text{#2}} & \footnotesize{#4}
    \end{tabular*}
    \vspace{-7mm}
}

\newcommand{\resumeEduOld}[5]{
\vspace{0mm}\item[]
    \begin{tabular*}{\textwidth}[t]{l@{\extracolsep{\fill}}r}
        \hspace{-4.3mm} \textbf{#1} & \footnotesize{#3}\vspace{-0.0mm} \\
        \hspace{-4.3mm} \small{#2} & \footnotesize{#4}\vspace{-0.0mm} \\ 
        \multicolumn{2}{l}{ \hspace{-4.3mm} \footnotesize{#5}}
    \end{tabular*}
    \vspace{-3.2mm}
}

\newcommand{\resumeEduNew}[2]{
\vspace{0mm}\item[]
    \begin{tabular*}{\textwidth}[t]{l@{\extracolsep{\fill}}r}
        \hspace{-4.3mm} \small\textbf{#1} & \footnotesize{#2}\vspace{-0.5mm} 
    \end{tabular*}
    \vspace{-3.2mm}
}

\newcommand{\resumeAchieve}[3]
{
\hspace{-3.1mm}\textbf{ #1} & {#2} & \hspace{3mm}\footnotesize{#3}
\vspace{0mm}\\
}


\renewcommand{\labelitemi}{$\vcenter{\hbox{\tiny$\bullet$}}$}

\newcommand{\resumeSubHeadingListStart}{\begin{itemize}[leftmargin=*,labelsep=0mm,itemsep=-2.5mm, topsep=1mm, partopsep=1mm, parsep=1mm]}

\newcommand{\resumeItemListStart}{\begin{justify}\begin{itemize}[leftmargin=1.5ex, rightmargin=2ex, noitemsep,labelsep=1.2mm,itemsep=1mm]\small}

\newcommand{\resumeSubHeadingListEnd}{\end{itemize}\vspace{-2mm}}
\newcommand{\resumeItemListEnd}{\end{itemize}\end{justify}}%\vspace{-1.5mm}}

\renewcommand{\arraystretch}{1}

\newcolumntype{L}{>{\raggedright\arraybackslash}X}%
\newcolumntype{R}{>{\raggedleft\arraybackslash}X}%
\newcolumntype{C}{>{\centering\arraybackslash}X}%

% Macro to create a dynamic table with specific column alignments without using the table environment
% #1 = Number of columns
% #2 = Table content
\newcommand{\InlineAlignedTable}[2]{%
    \centering
    \begin{tabularx}{\textwidth}{@{} L *{\numexpr#1-2\relax}{C} R @{}}
    #2
    \end{tabularx}
}

% Define helper commands for full-cell links
\newcommand{\FullCellLinkL}[2]{\href{#1}{\makebox[\linewidth][l]{#2}}}
\newcommand{\FullCellLinkC}[2]{\href{#1}{\makebox[\linewidth][c]{#2}}}
\newcommand{\FullCellLinkR}[2]{\href{#1}{\makebox[\linewidth][r]{#2}}}


%---- End of Packages and Functions ------



%-------------------------------------------
%%%%%%  CV STARTS HERE  %%%%%%%%%%%

\newcommand{\name}{Shantanu Nitin Ghodgaonkar} % Your Name
\newcommand{\phone}{+1 (929) 922-0614} % Your Phone Number
\newcommand{\emaila}{shantanu.ghodgaonkar@gmail.com} %Email 1
\newcommand{\github}{shantanu-ghodgaonkar} %Github



\begin{document}
\fontfamily{cmr}\selectfont


%----------HEADING-----------------
\begin{center}
    \LARGE{\textbf{\name}}
\end{center}
\vspace{-7.5mm}
\begin{justify}
    \href{https://shantanu-ghodgaonkar.github.io/}{shantanu-ghodgaonkar.github.io} | \href{https://www.linkedin.com/in/s-n-g/}{linkedin.com/in/s-n-g} | \href{mailto:\emaila}{\emaila} | \href{tel:+19299220614}{\phone}
\end{justify}
\vspace{-6.5mm}


%-----------EDUCATION-----------------

% \vspace{-2.5mm}
\section{Education}

\resumeSubHeadingListStart
\resumeEduNew
{New York University\textnormal{, Tandon School of Engineering}} 
{M.Sc. Mechatronics, Robotics and Automation Engineering} 
\resumeSubHeadingListEnd
\vspace{-2.5mm}

\resumeSubHeadingListStart
\resumeEduNew
{Visvesvaraya Technological University\textnormal{, Bangalore Institute of Technology}} 
{B.E. Electronics and Instrumentation Engineering} 
\resumeSubHeadingListEnd
\vspace{-4mm}



%-----------EXPERIENCE-----------------
% \vspace{-2.5mm}
\section{Experience}
\resumeSubHeadingListStart

\resumeExp
{Robotics Engineer\textnormal{ | BotCrew LLC | TX, USA}}
{\textit{Jun 2025 – Present}}
\resumeItemListStart
\item[$\bullet$] Leading full-stack software development of a \textbf{7-DoF robotic arm} integrated atop a \textbf{mobile robot}.
\item[$\bullet$] Built a \textbf{modular, platform-independent C++ framework} for control and planning across \textbf{ROS 2, Fast DDS, and RTOS}.
\item[$\bullet$] Designed a \textbf{collision-aware inverse kinematics (IK)} solver converging in \textbf{70–80 ms} to produce stable goal states.
\item[$\bullet$] Developed a joint-space \textbf{Model Predictive Controller (MPC)} running at \textbf{20 Hz} with position and velocity limits.
\item[$\bullet$] Implemented a \textbf{motion-planning pipeline} using \textbf{OMPL (RRT-Connect)} with self- and environment-collision checks.
\item[$\bullet$] Built a \textbf{PCD-to-SDF} module enabling \textbf{environment-aware motion planning from 3D sensor data}.
\item[$\bullet$] Programmed \textbf{low-level communication and actuator-control loops} for \textbf{hydraulic actuation} and future servo integration.
\item[$\bullet$] Developed a \textbf{Zero-G Controller} mode within the same ecosystem for human guidance control input.
\item[$\bullet$] Validating control and planning algorithms in \textbf{CoppeliaSim, MuJoCo and Isaac Lab}, targeting \textbf{Jetson Thor} deployment.
\resumeItemListEnd

\resumeExp
{Adjunct Professor\textnormal{ | New York University | NY, USA}}
{\textit{Jun 2024 - Jun 2025}}
\resumeItemListStart
\item[$\bullet$] Taught \textbf{control systems} using \textbf{C++, Python, MATLAB, and Simulink}, focusing on \textbf{PID, LQR, and MPC}.
\item[$\bullet$] Instructed students on \textbf{Standard Operating Procedures} for lab equipment including \textbf{oscilloscopes and function generators}.
\item[$\bullet$] Demonstrated the use of the \textbf{Allen Bradley PLC} to students, getting them ready for \textbf{real-world challenges}.
\item[$\bullet$] Designed and troubleshot \textbf{motion planning} for a \textbf{mobile hexapod} using \textbf{ROS Humble} and real-world testing.
\item[$\bullet$] Monitored and optimized \textbf{robotic performance} by debugging \textbf{simulation-to-hardware discrepancies} in \textbf{MuJoCo}.
\item[$\bullet$] Utilized \textbf{Linux, Bash, and Git} for \textbf{version control, system integration, and testing} for the \textbf{hexapod}.
\item[$\bullet$] Documented \textbf{control algorithms, system configurations, and troubleshooting steps} for improved \textbf{system reliability}.
\item[$\bullet$] Guided students in \textbf{debugging robotics issues}, adapting explanations to different skill levels for better learning.
\item[$\bullet$] Researching and testing applicable \textbf{Reinforcement Learning} strategies for \textbf{Vision-based Control} of the mobile hexapod.
% \item[$\bullet$] Currently leading the implementation of \textbf{vision-based control algorithms} for the hexapod.
% \item[$\bullet$] Passionate about \textbf{robotics and automation}, driving improvements in \textbf{hands-on learning} and \textbf{system deployment}.
\resumeItemListEnd

\resumeExp
{Software Engineer\textnormal{ | Bosch Global Software Technologies | Bengaluru, India}}
{\textit{Sep 2021 - Jul 2023}}
\resumeItemListStart
\item[$\bullet$] Built and maintained \textbf{diagnostic tools} for \textbf{ODX data processing}, ensuring seamless integration into \textbf{automotive systems}.
\item[$\bullet$] Identified and resolved \textbf{software issues} in \textbf{diagnostic tools}, enhancing performance and reducing \textbf{debugging time} for \textbf{INEOS}.
\item[$\bullet$] Documented \textbf{system errors} and submitted \textbf{Jira tickets} for issues requiring \textbf{senior engineering team support}.
\item[$\bullet$] \textbf{Collaborated} with engineering teams to improve \textbf{software stability}, cutting \textbf{development time} by \textbf{70\%} with \textbf{automation}.
\item[$\bullet$] Developed \textbf{Python scripts} to \textbf{automate repetitive tasks}, significantly improving \textbf{personal efficiency} by \textbf{40\%}.
\item[$\bullet$] Used \textbf{Git} and \textbf{Subversion} for \textbf{version control} and set up a \textbf{Jenkins CI/CD pipeline} for \textbf{automated deployment}.
\item[$\bullet$] Updated \textbf{documentation} on \textbf{sys-op, troubleshooting, and deployment process}, ensuring \textbf{clear knowledge transfer}.
\item[$\bullet$] Trained \textbf{client technicians} and \textbf{internal teams} on \textbf{new software workflows}, improving \textbf{operational efficiency}.
\item[$\bullet$] Conducted \textbf{on-site debugging} and \textbf{system validation}, \textbf{deploying} and \textbf{integrating tools} for \textbf{various clients}.
\item[$\bullet$] Worked in \textbf{Agile sprints} using \textbf{Jira}, maintaining a \textbf{90\% on-time delivery rate} while managing \textbf{multiple automotive projects}.
\item[$\bullet$] Took on additional responsibilities, including \textbf{client-facing roles}, \textbf{training new joiners}, and \textbf{conducting interviews}.
\item[$\bullet$] Earned \textbf{recognition and rewards} from \textbf{management} for contributions in improving \textbf{team performance} and \textbf{client satisfaction}.
\resumeItemListEnd




% \resumeExp
% {Diagnostic Content Engineering Intern\textnormal{ | Bosch Global Software Technologies | Bengaluru, India}}
% {\textit{Mar 2021 - Jun 2021}}
% \resumeItemListStart
% \item[$\bullet$] Developed \textbf{OTX screens} for \textbf{ECU simulation} using \textbf{HTML, CSS, and JavaScript} to support \textbf{diagnostic workflows}.
% \item[$\bullet$] Integrated and validated \textbf{OTX screens} for \textbf{vehicle testing} in \textbf{diagnostic systems}.
% \item[$\bullet$] Collaborated with \textbf{cross-functional teams} to improve \textbf{ECU simulation accuracy} and streamline \textbf{workflows}.
% \item[$\bullet$] Trained in \textbf{AUTOSAR, UDS and CAN} protocols and obtained \textbf{Road vehicles – Functional safety (ISO 26262)} certification.
% \resumeItemListEnd


% \resumeExp
% {Summer Engineering Intern\textnormal{ | FluxGen Sustainable Technologies | Bengaluru, India}}
% {\textit{Jul 2020 – Sep 2020}}
% \resumeItemListStart
% \item[$\bullet$] Developed a \textbf{wireless temperature and humidity monitoring system} using the \textbf{ESP32 Wi-Fi \& Bluetooth module}.
% \item[$\bullet$] Created an \textbf{Android app} to display \textbf{sensor data} in \textbf{real-time}, improving \textbf{monitoring accessibility}.
% \item[$\bullet$] Developed an \textbf{ad hoc wireless system} connecting \textbf{patients and doctors}, improving care during the \textbf{COVID-19 pandemic}.
% \item[$\bullet$] Designed \textbf{wearable devices} with \textbf{ESP32} $\mu$C and \textbf{LoRaWAN (RFM95)} for \textbf{wireless communication}.
% \item[$\bullet$] Integrated \textbf{sensors (MCP9808, MAX30102)} to track \textbf{body temperature, heart rate, and SpO2 levels} in \textbf{real-time}.
% \item[$\bullet$] Implemented a \textbf{meshed network} using the \textbf{ClusterDuck Protocol} for \textbf{real-time data collection} between \textbf{wearable devices}.
% \item[$\bullet$] Developed a \textbf{web-based user interface} using \textbf{HTML and JavaScript}, for \textbf{monitoring data} from \textbf{wearable devices}.
% \resumeItemListEnd


\resumeSubHeadingListEnd




% %-----------Projects-----------------

% % \vspace{-2.5mm}
% \section{Projects}
% \resumeSubHeadingListStart
% \resumeExp
% {Obstacle-Aware Control of 2D Quadrotor using PPO in Custom Reinforcement Learning Environment | \href{https://github.com/shantanu-ghodgaonkar/NYUMSMRFALL24/blob/main/Reinforcement Learning and Optimal Control for Robotics (ROB-GY 6323)/Project 2/project2.ipynb}{\underline{GitHub}}}{Dec 2024} \\
% \resumeExp
% {Deep Q-Learning to solve Linear Inverted Pendulum problem using PyTorch | \href{https://github.com/shantanu-ghodgaonkar/NYUMSMRFALL24/blob/main/Reinforcement Learning and Optimal Control for Robotics (ROB-GY 6323)/Homework/Series 4/Series 4 - Q-learning.ipynb}{\underline{GitHub}}}{Nov 2024} \\
% \resumeExp
% {Comparison between Value and Policy Iteration over Grid World | \href{https://github.com/shantanu-ghodgaonkar/NYUMSMRFALL24/blob/main/Reinforcement Learning and Optimal Control for Robotics (ROB-GY 6323)/Homework/Series 3/Exercise2.ipynb}{\underline{GitHub}}}{Oct 2024} \\
% \resumeExp
% {Trajectory Optimization and MPC for a 2D Quadrotor | \href{https://github.com/shantanu-ghodgaonkar/NYUMSMRFALL24/tree/main/Reinforcement Learning and Optimal Control for Robotics (ROB-GY 6323)/Project 1/Shantanu_Ghodgaonkar_ROBGY6323_Project1_Report.pdf}{\underline{GitHub}}}{Sep 2024} \\
% \resumeExp
% {Hexapod Robot Development | \href{https://github.com/shantanu-ghodgaonkar/hexapod}{\underline{GitHub}}}{Sep 2024 – Present} \\
% \resumeExp
% {Maze Solving Robot using Raspberry Pi and Arduino | \href{https://github.com/shantanu-ghodgaonkar/rPiUnoMazeSolver}{\underline{GitHub}}}{May 2024} \\
% \resumeExp
% {Vision and IMU Fusion with Unscented Kalman Filter | \href{https://github.com/shantanu-ghodgaonkar/Proj3-ROBGY6213}{\underline{GitHub}}}{April 2024} \\
% \resumeSubHeadingListEnd


% \vspace{-3.5mm}


%-----------SKILLS-----------------
\vspace{-4.5mm}
\section{Technical Skills}
% \vspace{0.2mm}

% \small{\begin{tabular*}{\textwidth}[t]{p{1\textwidth}}% p{0.5\textwidth}}
% % Robotics fundamentals, including geometry, linear algebra, multivariate calculus, kinematics, dynamics, probability, and statistics 
%     \hspace{-3.1mm}{\textbf{ Control Systems:} PID Control, LQR Control, Model Predictive Control, Numerical Optimization, PLC Ladder Logic} \vspace{1mm}\\
%     % \hspace{-3.1mm}{\textbf{ Robotics Fundamentals:} 3D and Epipolar Geometry, Linear Algebra, Kinematics, Dynamics, Probability, Statistics} \vspace{1mm}\\ 
%     \hspace{-3.1mm}{\textbf{ Robotics Platforms:} Universal Robotics UR16, MuJoCo, Gazebo, CoppeliaSim, Nvidia Jetson, Raspberry Pi, Arduino, ESP32} \vspace{1mm}\\ 
%     \hspace{-3.1mm}{\textbf{ Programming Languages:} Python, C++, C, Java, HTML, CSS, JavaScript, XML, SQL, Linux Bash} \vspace{1mm}\\ 
%     \hspace{-3.1mm}{\textbf{ Frameworks \& Libraries:} PyTorch, ROS Humble, OpenCV, SciPy, Pinocchio, Simulink, MATLAB Robotics Toolbox} \vspace{1mm}\\ 
%     \hspace{-3.1mm}{\textbf{ Communication Protocols:} UART, USB, I2C, SPI, ClusterDuck Protocol, Dynamixel Protocol 2.0, BLE, WiFi, MQTT, CAN} \vspace{1mm}\\ 
%     \hspace{-3.1mm}{\textbf{ Tools \& Others:} Git, Subversion, Docker, Jira, LabVIEW, LPKF CircuitPro, Altium, Eagle, KiCad, Overleaf} \vspace{1mm}
% \end{tabular*}}



% \section{Technical Skills}
\small{\begin{tabular*}{\textwidth}[t]{p{1\textwidth}}
    \hspace{-3.1mm}{\textbf{ Control Systems:} PID, LQR, Numerical Optimization, ProxQP, OSQP, Model Predictive Control} \vspace{1mm}\\
    \hspace{-3.1mm}{\textbf{ Robotics \& Motion Planning:} Pinocchio, OMPL, nvblox, CoppeliaSim, MuJoCo, ROS 2 Humble, CycloneDDS} \vspace{1mm}\\
    \hspace{-3.1mm}{\textbf{ Programming \& Tools:} C++, Python, C, Embedded C, CUDA, CMake, Docker, Git, Linux Bash} \vspace{1mm}\\
    \hspace{-3.1mm}{\textbf{ Simulation \& Hardware:} Nvidia Jetson, Raspberry Pi, Arduino, ESP32, STM32 (Moteus X1)} \vspace{1mm}\\
    \hspace{-3.1mm}{\textbf{ Libraries \& Frameworks:} PyTorch, OpenCV, SciPy, MATLAB/Simulink, MATLAB Robotics Toolbox} \vspace{1mm}\\
    \hspace{-3.1mm}{\textbf{ Communication Protocols:} UART, USB, I2C, SPI, CAN, CAN-FD, BLE, WiFi, MQTT, Dynamixel Protocol 2.0} \vspace{1mm} \\
    \hspace{-3.1mm}{\textbf{ Tools \& Others:} Git, Subversion, Docker, Jira, LabVIEW, LPKF CircuitPro, Altium, Eagle, KiCad, Overleaf} \vspace{1mm}
\end{tabular*}}




\end{document}
